\documentclass[11pt]{amsart}
\usepackage{amsmath,amssymb,amsthm}
\usepackage[margin=1.15in]{geometry}
\usepackage{booktabs}
\usepackage[hidelinks]{hyperref}
\hypersetup{pdftitle={The superbase of a uniform spanning tree on a torus},pdfauthor={Yitzchak Shmalo}}
\newtheorem{theorem}{Theorem}[section]
\newtheorem{lemma}[theorem]{Lemma}
\theoremstyle{definition}
\newtheorem{definition}[theorem]{Definition}
\theoremstyle{remark}
\newtheorem{remark}[theorem]{Remark}
\newcommand{\Z}{\mathbb{Z}}
\newcommand{\R}{\mathbb{R}}
\newcommand{\C}{\mathbb{C}}
\newcommand{\Prob}{\mathbb{P}}
\newcommand{\doi}[1]{\href{https://doi.org/#1}{doi:\nolinkurl{#1}}}
\newcommand{\arxiv}[1]{\href{https://arxiv.org/abs/#1}{arXiv:\nolinkurl{#1}}}
\setlength{\emergencystretch}{2.5em}

\title[The superbase of a uniform spanning tree on a torus]{The superbase of a uniform spanning tree on a torus}
\author{Yitzchak Shmalo}
\address{Einstein Institute of Mathematics, The Hebrew University of Jerusalem, Jerusalem, Israel}
\email{yspennstate@gmail.com}
\date{November 2026}
\subjclass[2020]{60D05, 82B20, 05C05, 05C81, 65C05}
\keywords{uniform spanning tree, torus, homology, conformal invariance, Wilson's algorithm, superbase}

\raggedbottom
\begin{document}
\begin{abstract}
Lyons asked for the law of the homology classes carried by the dual of a uniform spanning tree on a fine grid on a
torus, and whether it has a conformally invariant limit. The nonzero classes of the fundamental cycles of the tree are
those of the dual graph: a superbase of the period lattice when the core of the dual is a theta graph, and a basis
when it is a figure-eight, which we record separately. We sample $10^6$ trees on each of $80$ tori with Wilson's
algorithm. Over three meshes at three moduli the leading probabilities move by at most $1.3\cdot10^{-3}$, the
figure-eight frequency falls by a factor of about $3.5$ per halving of the mesh, and tori placed at an angle to the grid
agree with the equivalent axis-parallel tori, which supports a conformally invariant limit. We give finest-mesh
estimates at four moduli and show that, on these meshes, the weight of a superbase is not an exponential of its
energy.
\end{abstract}
\maketitle

\section{Introduction}

Let $\Sigma_\tau=\C/(\Z+\tau\Z)$ with $\operatorname{Im}\tau>0$, and let $G$ be a square grid on $\Sigma_\tau$; for
rational moduli such grids exist at arbitrarily fine mesh, and every modulus is a limit of rational ones. A spanning tree
$T$ of $G$ determines a dual graph $D(T)$ on $\Sigma_\tau$, whose vertices are the faces of $G$ and whose edges are the
dual edges $e^*$ of the edges $e\notin T$. Lyons observed that $D(T)$ contains only noncontractible cycles and asked
for the law of the homology basis it gives when $T$ is uniform~\cite[Section~12]{Ly03}: ``Does it have a limit as the
mesh of the graph tends to 0?'' He expected the limit to exist and to be conformally invariant. Schramm included the
question in his collection of problems on conformally invariant scaling limits, as a problem proposed by
Lyons~\cite[Problem~6.2]{Sch07}. In his formulation $D(T)$ contains exactly three edge-simple closed paths, and the task
is to ``[d]etermine the distribution of the triple of homotopy classes containing these three closed paths'' as a
function of the geometry of the torus. On a torus free homotopy classes of closed curves are homology classes, so up to
orientation the triple is a set of three classes in $H_1(\Sigma_\tau;\Z)\cong\Z^2$.

On graphs embedded in surfaces of positive genus the closest exact results concern cycle-rooted spanning forests. Forman
expressed the determinant of the Laplacian twisted by a connection as a sum over such forests~\cite{For93}. Kenyon used
the Laplacian twisted by a flat line bundle to compute the law of the homology classes of their cycles, with an explicit
calculation on the square torus~\cite{Ken11}, and Kassel and Kenyon proved conformally invariant scaling limits for the
resulting measures on noncontractible multicurves~\cite{KK17}. Kenyon and Wilson gave a method for computing the
probabilities of the connection topologies of a uniform spanning tree on a graph embedded in a surface, and applied it
to loop-erased random walk on planar graphs~\cite{KW15}. For dimers on tori, Kenyon, Sun and Wilson found the
asymptotics of the partition functions and the distribution of the height change around the torus~\cite{KSW16}.
Dub\'edat identified the compactified free field as the limit of the toroidal dimer height~\cite{Dub15}. Dub\'edat and
Gheissari showed that the height periods of toroidal Temperleyan dimers have a universal discrete Gaussian
law~\cite{DG15}. Berestycki, Laslier and Ray developed the Temperleyan correspondence on Riemann surfaces and proved
conformally invariant limits for the associated forests~\cite{BLR25,BLR24}. Lam, Lo and Yuen showed that the dual
complement of a homological quasi-tree on a discrete surface is again a homological quasi-tree; on a torus the dual of
a spanning tree is one of rank two~\cite{LLY26}. These results concern other observables: disjoint cycles of forests,
height periods, or connection events. An explicit joint law of the superbase is not known to us.

We study the question numerically. Section~\ref{sec:struct} proves that the fundamental cycles of $T$ carry the same
classes as $D(T)$, so that sampling $T$ is enough. Section~\ref{sec:method} describes the sampler, Wilson's
algorithm~\cite{Wil96}, and checks it against exact enumeration. Section~\ref{sec:results} reports the estimates. The
main observations are these.
\begin{enumerate}
\item The frequency of a figure-eight core of $D(T)$ falls by a factor of $3.2$ to $3.9$ each time the mesh of the
square torus is halved (Table~\ref{tab:eight}), consistent with a limit in which the nonzero classes always form a superbase.
\item Over three meshes at $\tau=i$, $\tau=0.5+i$ and $\tau$ near $e^{i\pi/3}$ the leading probabilities move by at most
$1.3\cdot10^{-3}$ (Table~\ref{tab:ladder}).
\item Tori placed at an angle to the grid give the law of the axis-parallel torus of the same conformal class within
sampling error (Table~\ref{tab:place}).
\item On the square torus each of the two superbases $\{1,i,1\pm i\}$ has frequency $0.475$ at the finest mesh; near
the hexagonal torus the equilateral superbase has frequency $0.728$ (Table~\ref{tab:limit}).
\item The weight of a superbase is close to $\exp(-0.9\,Q)$ for the energy $Q$ of~\eqref{eq:Q} near the square and
hexagonal tori, but not at $\tau=2i$; polynomials in two shape invariants fit it on a grid of $49$ moduli, and their
extrapolation to longer tori is poorly determined (Table~\ref{tab:fit}).
\end{enumerate}
A formula for the limit law is not found here. The tables give the values a candidate formula has to reproduce.

\section{The classes of the dual graph}\label{sec:struct}

Let $\Lambda\subset\Z^2$ be the lattice generated by $(n_1,0)$ and $(s,n_2)$, with $n_1,n_2\ge3$ and $0\le s<n_1$. We
take $G=\Z^2/\Lambda$, embedded in the torus $\Sigma=\R^2/\Lambda$ with unit square faces; the condition $n_1,n_2\ge3$
makes $G$ a simple $4$-regular graph. Dividing by $n_1$ maps $\Sigma$ to $\Sigma_\tau$ with $\tau=(s+in_2)/n_1$. A
closed curve on $\Sigma$ lifts to a path in $\R^2$ with displacement $a(n_1,0)+b(s,n_2)$ for integers $a,b$. We call
$(a,b)$, or the period $a+b\tau$, its class, and identify $H_1(\Sigma;\Z)$ with $\Z^2$ in this way. The algebraic
intersection number of the classes $(a,b)$ and $(c,d)$ is $ad-bc$.

\begin{definition}\label{def:superbase}
A superbase of $\Z^2$ is a set $\{v_1,v_2,v_3\}$ of classes, each taken up to sign, whose signs can be chosen so that
$v_1+v_2+v_3=0$, and any two of which form a basis. Its energy is
\begin{equation}\label{eq:Q}
  Q(S)=\frac{|v_1|^2+|v_2|^2+|v_3|^2}{\operatorname{Im}\tau},
\end{equation}
where $|v|$ is the length of the period $v=a+b\tau$.
\end{definition}

The three periods of a superbase are the sides of a triangle of area $\Delta=\operatorname{Im}\tau/2$. For a triangle
with angles $A,B,C$ one has $|v_1|^2+|v_2|^2+|v_3|^2=4\Delta(\cot A+\cot B+\cot C)$, so $Q(S)$ is twice the sum of the
cotangents of its angles and depends only on the shape of the triangle.

For $e\notin T$ let $\gamma_e$ be the cycle formed by $e$ and the path in $T$ joining its endpoints, and let $h(e)$ be
its class up to sign.

\begin{lemma}\label{lem:struct}
Let $T$ be a spanning tree of $G$ and $D=D(T)$ its dual graph.
\begin{enumerate}
\item $D$ is connected, its cyclomatic number is $2$, and the inclusion induces an isomorphism $H_1(D;\Z)\to
H_1(\Sigma;\Z)$.
\item The core of $D$, obtained by repeatedly removing vertices of degree one, is a subdivided theta graph or a
subdivided figure-eight.
\item Suppose the core is a theta with arcs $\alpha_1,\alpha_2,\alpha_3$, and let $p_i$ be the class of the cycle
$\alpha_j\cup\alpha_k$, $\{i,j,k\}=\{1,2,3\}$. Then $\{p_1,p_2,p_3\}$ is a superbase, $h(e)=\pm p_i$ when $e^*$ lies
on $\alpha_i$, and $h(e)=0$ when $e^*$ is not on the core.
\item Suppose the core is a figure-eight whose loops have classes $d_1,d_2$. Then $\{d_1,d_2\}$ is a basis,
$h(e)=\pm d_2$ when $e^*$ lies on the first loop, $h(e)=\pm d_1$ when it lies on the second, and $h(e)=0$ otherwise.
\end{enumerate}
\end{lemma}

\begin{proof}
(1) A closed regular neighbourhood $N$ of $T$ in $\Sigma$ is a closed disc. The complement of its interior deformation
retracts onto $D$: each face of $G$ minus $N$ retracts onto its centre and the adjacent dual edges, and the part of a
non-tree edge outside $N$ retracts onto the point where it meets its dual edge. The complement of an open disc in a
torus is a torus with one hole, and its first homology maps isomorphically onto that of the torus, because the
boundary circle is a commutator in the fundamental group. So $H_1(D)\to H_1(\Sigma)$ is an isomorphism onto $\Z^2$,
$D$ is connected, and its cyclomatic number is $2$.

(2) Removing a vertex of degree one does not change the homotopy type. So the core $K$ is connected, has no vertex of
degree one, and has cyclomatic number $2$. After suppressing the vertices of degree two, every vertex of $K$ has
degree at least three and $|E|-|V|=1$, hence $3|V|\le2|E|=2|V|+2$ and $|V|\le2$. With two vertices $K$ is a theta
graph or a dumbbell, two loops joined by an edge; with one vertex it is a figure-eight. The loops of a dumbbell are
disjoint closed curves, so their intersection number is $0$. By (1) their classes form a basis, so their intersection
number is $\pm1$. Hence $K$ is not a dumbbell.

(3) The cycle $\gamma_e$ is a simple closed curve that meets $D$ in exactly one point, where $e$ crosses $e^*$
transversally: the other edges of $\gamma_e$ are in $T$, so their duals are not in $D$, and the vertices of $G$ are
not on $D$. Hence for every cycle $z$ of $D$ the intersection number of $h(e)$ with the class of $z$ is, up to sign,
the signed number of times $z$ traverses $e^*$. If $e^*$ is not on the core, no cycle of $D$ uses it. By (1) the
classes of cycles of $D$ exhaust $\Z^2$, and the intersection form is nondegenerate, so $h(e)=0$. For suitable
orientations $p_1\pm p_2\pm p_3=0$, and any two of the $p_i$ form a basis by (1), so $\{p_1,p_2,p_3\}$ is a superbase.
If $e^*$ lies on $\alpha_1$, the cycle $\alpha_2\cup\alpha_3$ avoids $e^*$ and the cycle $\alpha_1\cup\alpha_2$
traverses it once. So $h(e)$ has intersection number $0$ with $p_1$ and $\pm1$ with $p_3$. As $p_1$ is primitive, the
first gives $h(e)=\lambda p_1$ for an integer $\lambda$, and the second gives $\lambda=\pm1$.

(4) By (1) the loop classes $d_1,d_2$ form a basis, so $d_1\cdot d_2=\pm1$ and both are primitive. If $e^*$ lies on the
first loop, the second loop avoids $e^*$ and the first traverses it once. So $h(e)$ has intersection number $0$ with
$d_2$ and $\pm1$ with $d_1$. The first gives $h(e)=\lambda d_2$, and then $\lambda d_2\cdot d_1=\pm1$ gives
$\lambda=\pm1$. The second loop is treated in the same way, and $h(e)=0$ off the core as in (3).
\end{proof}

\begin{remark}
Part (1) is the case $g=1$, $k=0$ of the correspondence of Lam, Lo and Yuen between a homological quasi-tree and the dual
complement of another~\cite{LLY26}; the argument above is the elementary one for a torus.
\end{remark}

So the set of nonzero classes $h(e)$, $e\notin T$, is the superbase $\{p_1,p_2,p_3\}$ in the theta case and the basis
$\{d_1,d_2\}$ in the figure-eight case. In the theta case the edge-simple closed paths of $D$ are the three cycles
$\alpha_j\cup\alpha_k$, as in Schramm's formulation. In the figure-eight case there are four up to reversal: the two
loops, of classes $d_1$ and $d_2$, and the two closed trails that run through both loops in either relative
orientation, of classes $d_1+d_2$ and $d_1-d_2$. For example, on the $3\times3$ torus let $T$ consist of all horizontal
edges that do not wrap around the torus and the two vertical edges of column $0$ that do not wrap. Then the core of
$D(T)$ is one horizontal and one vertical loop meeting at a vertex, with classes $(1,0)$ and $(0,1)$. So at a finite
mesh the triple of Schramm's formulation is defined only on the theta event, and we record figure-eight trees
separately. They are a sizeable part of the small tori: $2{,}592$ of the $11{,}664$ spanning trees of the $3\times3$
torus, $1{,}674$ of the $12{,}321$ of the torus $(3,3,1)$ and $63{,}900$ of the $367{,}500$ of the torus $(4,3,0)$.
Section~\ref{sec:results} shows that their frequency falls quickly with the mesh.

\section{Method}\label{sec:method}

\subsection*{Sampling}
A uniform spanning tree of $G$ is sampled by Wilson's algorithm~\cite{Wil96}. Starting from a root, each vertex not
yet in the tree starts a simple random walk, which is stopped when it hits the tree, and its loop erasure is added to
the tree. The tree is then lifted to $\Z^2$, giving each vertex a point of $\Z^2$ in its class modulo $\Lambda$. The
CPU program follows the parent pointers from a vertex to the first vertex already lifted and unwinds the chain; the GPU
program lifts each loop-erased path when it is attached. For a non-tree edge from $u$ to $v$ in direction
$(\Delta_x,\Delta_y)$ the displacement of $\gamma_e$ is $(X,Y)=\mathrm{lift}(u)+(\Delta_x,\Delta_y)-\mathrm{lift}(v)$,
and its class is $b=Y/n_2$, $a=(X-bs)/n_1$. By Lemma~\ref{lem:struct} the distinct nonzero classes are three in the
theta case and two in the figure-eight case.

Two implementations were used: a CPU program compiled with Numba, and a program for an Apple GPU in which each thread
samples one tree by the same algorithm. The CPU program draws tree $k$ of a run with seed $\sigma$ from Numba's Mersenne
Twister reseeded with $1000003\,\sigma+k$. The GPU program gives each thread a PCG generator initialised from the seed
and the index of its tree. On a Mac mini the GPU program samples $10^6$ trees of the $160\times160$ torus in $206$
seconds.

\subsection*{Validation}
On three small tori we listed every spanning tree, by testing every set of $N-1$ of the $2N$ edges with a union-find
structure, and computed the exact law of the set of nonzero classes with a separate program. The numbers of spanning
trees found agree with the matrix-tree theorem. Table~\ref{tab:exact} compares the exact laws with $10^6$ draws of
each sampler, with fixed seeds. No draw produced a class set of probability zero.

\begin{table}[tbp]
\centering
\small
\begin{tabular}{lcccc}
\toprule
torus $(n_1,n_2,s)$ & spanning trees & class sets & GPU, $\chi^2$/d.f. & CPU, $\chi^2$/d.f.\\\midrule
$(3,3,0)$ & $11{,}664$ & $11$ & $9.4/10$ & $10.8/10$\\
$(3,3,1)$ & $12{,}321$ & $9$ & $8.8/8$ & $2.5/8$\\
$(4,3,0)$ & $367{,}500$ & $13$ & $11.6/12$ & $5.0/12$\\\bottomrule
\end{tabular}
\medskip
\caption{Samples against the exact law of the set of nonzero classes of fundamental cycles, $10^6$ trees for each
sampler and torus. The class sets include the two-class sets of figure-eight trees.}
\label{tab:exact}
\end{table}

\subsection*{Sample sizes and errors}
Every frequency reported below is from $10^6$ trees. A superbase is written by its periods. Standard errors are
binomial, about $5\cdot10^{-4}$ for the leading superbases. A difference between two tori is measured in combined
standard errors, the square root of the sum of the two squared standard errors.

\section{Results}\label{sec:results}

\subsection*{The figure-eight}
Table~\ref{tab:eight} gives the frequency of the figure-eight on the square torus. It falls by a factor between $3.2$
and $3.9$ for each halving of the mesh, and on the other two ladders of Table~\ref{tab:ladder} by similar factors.
This is consistent with a limit in which the core of $D(T)$ is a theta graph and the nonzero classes form a random
superbase.

\begin{table}[tbp]
\centering
\small
\begin{tabular}{lcccccc}
\toprule
$n$ & 8 & 16 & 32 & 64 & 128 & 256\\\midrule
figure-eight & $0.0223$ & $0.0058$ & $0.0016$ & $0.00044$ & $0.00013$ & $0.00004$\\\bottomrule
\end{tabular}
\medskip
\caption{Frequency of a figure-eight core on the $n\times n$ torus, $10^6$ trees each.}
\label{tab:eight}
\end{table}

\subsection*{Convergence}
Table~\ref{tab:ladder} follows the leading superbases over three meshes at three moduli. At $\tau=i$ and $\tau=0.5+i$
the largest change from the coarsest to the finest mesh is $1.22\cdot10^{-3}$, for $\{1,i,1+i\}$ at $\tau=i$, which is
$1.7$ combined standard errors. Near the hexagonal modulus the superbase $\{1,\tau,1+\tau\}$ moves by
$1.30\cdot10^{-3}$, or $3.2$ combined standard errors, from the coarsest to the finest mesh, and by at most $1.7$
combined standard errors between the two finer meshes. The table gives no bound on the discretization error, but the
changes are of the order of the sampling error.

\begin{table}[tbp]
\centering
\small
\begin{tabular}{llccc}
\toprule
$\tau$ & superbase & coarse & middle & fine\\\midrule
$i$ & $\{1,i,1+i\}$ & $0.47405$ & $0.47486$ & $0.47527$\\
 & $\{1,i,1-i\}$ & $0.47455$ & $0.47477$ & $0.47449$\\
 & the four with $Q=8$ & $0.01237$--$0.01256$ & $0.01239$--$0.01259$ & $0.01245$--$0.01258$\\
 & figure-eight & $0.00158$ & $0.00044$ & $0.00013$\\
$0.5+i$ & $\{1,\tau,1-\tau\}$ & $0.71419$ & $0.71420$ & $0.71371$\\
 & $\{1,\tau,1+\tau\}$ & $0.11833$ & $0.11886$ & $0.11921$\\
 & $\{1,1-\tau,2-\tau\}$ & $0.11850$ & $0.11877$ & $0.11936$\\
 & $\{\tau,1-\tau,1-2\tau\}$ & $0.04654$ & $0.04641$ & $0.04617$\\
 & figure-eight & $0.00099$ & $0.00029$ & $0.00009$\\
$0.5+\tfrac{13}{15}i$ & $\{1,\tau,1-\tau\}$ & $0.72731$ & $0.72769$ & $0.72756$\\
 & $\{1,\tau,1+\tau\}$ & $0.08953$ & $0.09071$ & $0.09083$\\
 & $\{1,1-\tau,2-\tau\}$ & $0.08988$ & $0.08995$ & $0.09065$\\
 & $\{\tau,1-\tau,1-2\tau\}$ & $0.09052$ & $0.09017$ & $0.08984$\\
 & figure-eight & $0.00184$ & $0.00053$ & $0.00014$\\\bottomrule
\end{tabular}
\medskip
\caption{Frequencies over three meshes, $10^6$ trees each: the $n\times n$ torus for $\tau=i$ ($n=32,64,128$), the
torus $(n,n,n/2)$ for $\tau=0.5+i$ ($n=40,80,160$), and the tori $(30,26,15)$, $(60,52,30)$, $(120,104,60)$ for
$\tau=0.5+\tfrac{13}{15}i$.}
\label{tab:ladder}
\end{table}

\subsection*{Placement on the grid}
Let $g$ be an element of $\mathrm{PSL}_2(\Z)$, acting on the modulus by a M\"obius transformation, so that the torus
of modulus $g\tau$ is conformally equivalent to $\Sigma_\tau$. If the limit is conformally invariant, the laws at
$\tau$ and at $g\tau$ agree once the classes are carried over by $g$, although the two tori sit differently on the
square grid. We realized $(\tau-1)/\tau=0.6+0.8i$ for $\tau=0.5+i$ as the torus $(200,160,120)$, and
$2i/(1+2i)=0.8+0.4i$ as $(300,120,240)$, and compared them with the axis-parallel tori $(160,160,80)$ and
$64\times128$. For each pair we formed the $2\times K$ table of theta-core trees by superbase, pooling the superbases
with fewer than $200$ trees in the two tori together into one cell, and computed Pearson's $\chi^2$.
Table~\ref{tab:place} gives the result. The first pair gives $\chi^2=10.09$ on $8$ degrees of freedom and the second
$7.24$ on $8$. In the first pair the four leading superbases agree within $0.34$ combined standard errors. In the
second pair the largest single difference, $2.1$ combined standard errors, is in one of the top
two superbases, which form a mirror pair; their sum differs by $1.27$ combined standard errors. The figure-eight
frequencies of the second pair differ, $362$ against $103$ per million, as expected from the mesh effect of
Table~\ref{tab:eight}: the torus $(300,120,240)$ has $36{,}000$ vertices and $64\times128$ has $8{,}192$.

\begin{table}[tbp]
\centering
\footnotesize
\setlength{\tabcolsep}{4pt}
\begin{tabular}{llll}
\toprule
pair & axis-parallel & at an angle & $\chi^2$/d.f.\\\midrule
$0.5+i$ and $0.6+0.8i$ & $0.7137$, $0.1192$, $0.1194$, $0.0462$ & $0.7137$, $0.1192$, $0.1193$, $0.0463$ & $10.09/8$\\
$2i$ and $0.8+0.4i$ & $0.4196$, $0.4207$; $0.0770$, $0.0768$ & $0.4211$, $0.4199$; $0.0765$, $0.0769$ & $7.24/8$\\\bottomrule
\end{tabular}
\medskip
\caption{Tori placed at an angle to the grid against the equivalent axis-parallel tori, $10^6$ trees each. The classes
are mapped by $\left(\begin{smallmatrix}0&-1\\1&1\end{smallmatrix}\right)$ and
$\left(\begin{smallmatrix}1&0\\1&1\end{smallmatrix}\right)$. The frequencies are listed in the order $\{1,\tau,1-\tau\}$, $\{1,\tau,1+\tau\}$, $\{1,1-\tau,2-\tau\}$,
$\{\tau,1-\tau,1-2\tau\}$ for the first pair, with $\tau=0.5+i$, and
$\{1,2i,1+2i\}$, $\{1,2i,1-2i\}$; $\{1,1+2i,2+2i\}$, $\{1,1-2i,2-2i\}$ of the axis-parallel torus.}
\label{tab:place}
\end{table}

\subsection*{Finest-mesh estimates at four moduli}
Table~\ref{tab:limit} gives the law on the finest mesh at four moduli. By Table~\ref{tab:ladder} we take these as
estimates of the conjectured limit to about three digits. On the $256\times256$ torus the six superbases listed carry
all but $79$ of $10^6$ trees: $43$ figure-eights and $36$ other theta cores. Mirror superbases have equal frequencies
within sampling error. The torus $(120,104,60)$ has modulus $0.5+\tfrac{13}{15}i$, at distance $6.4\cdot10^{-4}$ from
$e^{i\pi/3}$; its lattice has only the mirror symmetry, and its three superbases of the second level have energies
$5.772$, $5.772$ and $5.777$.

\begin{table}[tbp]
\centering
\small
\begin{tabular}{llcc}
\toprule
$\tau$ (torus) & superbase & $Q$ & frequency\\\midrule
$i$ ($256\times256$) & $\{1,i,1+i\}$, $\{1,i,1-i\}$ & $4$ & $0.4750$, $0.4749$\\
 & four, e.g.\ $\{1,1+i,2+i\}$ & $8$ & $0.0124$--$0.0126$\\
 & figure-eight & & $0.00004$\\
$0.5+\tfrac{13}{15}i$ ($(120,104,60)$) & $\{1,\tau,1-\tau\}$ & $3.464$ & $0.7276$\\
 & $\{1,\tau,1+\tau\}$, $\{1,1-\tau,2-\tau\}$ & $5.772$ & $0.0908$, $0.0907$\\
 & $\{\tau,1-\tau,1-2\tau\}$ & $5.777$ & $0.0898$\\
 & six, e.g.\ $\{1,1+\tau,2+\tau\}$ & $12.70$ & $0.00015$--$0.00018$\\
$0.5+i$ ($(160,160,80)$) & $\{1,\tau,1-\tau\}$ & $3.5$ & $0.7137$\\
 & $\{1,\tau,1+\tau\}$, $\{1,1-\tau,2-\tau\}$ & $5.5$ & $0.1192$, $0.1194$\\
 & $\{\tau,1-\tau,1-2\tau\}$ & $6.5$ & $0.0462$\\
$2i$ ($64\times128$) & $\{1,2i,1+2i\}$, $\{1,2i,1-2i\}$ & $5$ & $0.4196$, $0.4207$\\
 & $\{1,1+2i,2+2i\}$, $\{1,1-2i,2-2i\}$ & $7$ & $0.0770$, $0.0768$\\
 & $\{1,2+2i,3+2i\}$, $\{1,2-2i,3-2i\}$ & $11$ & $0.0026$, $0.0027$\\
 & figure-eight & & $0.00036$\\\bottomrule
\end{tabular}
\medskip
\caption{The law of the superbase on the finest mesh, $10^6$ trees each. Standard errors are about $5\cdot10^{-4}$
for the leading values and $10^{-4}$ near $0.01$.}
\label{tab:limit}
\end{table}

\subsection*{The weight of a superbase}
Suppose that the law conditioned on a theta core has a limit that is conformally invariant and invariant under
reflection. A superbase determines its lattice up to similarity, and reflection maps the torus of modulus $\tau$ to
that of modulus $-\bar\tau$, so the limiting probability of a superbase $S$ is then a function of the unoriented shape
of its triangle alone. At a finite mesh we write the frequencies as $\Prob_\tau(S)=w(\text{shape of }S)/Z(\tau)$, with
a positive weight $w$ and a normalisation $Z(\tau)$ for each torus, which also absorbs the figure-eight mass and the
superbases left out of a fit. The content of a formula is in the form of $w$. The shape of a triangle with angles
$A,B,C$ is fixed by $e_1=\cot A+\cot B+\cot C=Q/2$ and $e_3=\cot A\cot B\cot C$: since $\cot A\cot B+\cot B\cot
C+\cot C\cot A=1$ for every triangle, the three cotangents are the roots of $t^3-e_1t^2+t-e_3$.

The simplest guess is $w=\exp(-cQ)$. The ratio of the top superbase to the next level on the finest mesh, with the
symmetric superbases of a level pooled, gives $c=0.909$ on the square torus ($Q$ from $4$ to $8$), $0.903$ near the
hexagonal torus ($Q$ from $3.46$ to $5.77$) and $0.911$ from there to $Q=12.7$, and $0.894$ and $0.913$ at $0.5+i$
($Q$ from $3.5$ to $5.5$ and to $6.5$). At $2i$ it gives $0.849$. The standard errors are about $0.001$ on the square
torus and at $2i$, so the difference of $0.06$ between them is far outside sampling error. On these meshes a single
exponential of $Q$ holds to a few percent near the square and hexagonal tori and fails on the longer torus.

To see more of $w$ we sampled a grid of $49$ moduli: the tori $(100,n_2,s)$ with
$n_2\in\{100,110,125,150,175,200,250,300\}$ and $s\in\{0,10,\dots,50\}$, and the torus $(100,90,50)$. This gives $333$
cells with at least $200$ counts and $284$ distinct shapes. We fitted $\log w$ as a polynomial in $(e_1,e_3)$ without
constant term by weighted least squares on the log-frequencies, with $\log Z$ free for each modulus and the weight of a
cell the square root of its count. Table~\ref{tab:fit} shows the fit by degree. The fourth-degree fit leaves
$\chi^2=338.7$ on $270$ degrees of freedom ($p\approx0.003$), so it is close but not exact on the grid.

\begin{table}[tbp]
\centering
\small
\begin{tabular}{lcccc}
\toprule
degree & 1 & 2 & 3 & 4\\\midrule
shape coefficients & 2 & 5 & 9 & 14\\
residual degrees of freedom & 282 & 277 & 273 & 270\\
$\chi^2$ per degree of freedom & 94.03 & 11.67 & 2.54 & 1.25\\\bottomrule
\end{tabular}
\medskip
\caption{Fit of $\log w$ by a polynomial in $(e_1,e_3)$ on the grid of $49$ moduli, with one free $\log Z$ per
modulus.}
\label{tab:fit}
\end{table}

We tested the fourth-degree polynomial on six finer meshes outside the grid: the tori $256\times256$, $(160,160,80)$,
$(200,160,120)$, $(300,120,240)$, $64\times128$ and $(120,104,60)$. With the coefficients held fixed and only $\log Z$
refitted, $\chi^2$ per degree of freedom is $1.02$ over $28$ degrees of freedom, with no cell beyond $1.92$ standard
errors. A predictive statistic that includes the uncertainty of the fitted coefficients, the increase of the minimised
$\chi^2$ when the six tori join the fit, each with its own $\log Z$, is $27.78$ on $28$ degrees of freedom
($p=0.48$).

In aspect ratio the fit is poorly determined. We sampled $18$ longer tori, $(100,n_2,s)$ with $n_2\in\{350,400,500\}$.
With the coefficients of the fifth-degree polynomial fitted on the grid held fixed, $\chi^2$ per degree of freedom is
$2.17$ at $n_2=350$ and $24.7$ at $n_2=400$. Most of this is the uncertainty of the coefficients outside the range of
the grid: the predictive statistic is $45.45$ on $43$ degrees of freedom at $n_2=350$ ($p=0.37$) and $43.23$ on $45$ at
$n_2=400$ ($p=0.55$). Fitted with $n_2=350$ and $400$ added, the fifth- and sixth-degree polynomials give $\chi^2$ per
degree of freedom $101$ and $99$ at $n_2=500$ with fixed coefficients, and predictive statistics $95.16$ on $51$
($p=1.7\cdot10^{-4}$) and $55.01$ on $51$ ($p=0.33$). So a polynomial in $(e_1,e_3)$ describes $w$ well on the range of
shapes it is fitted on, and the data do not determine it beyond that range. Short closed forms do worse on the grid. We
tried combinations of $Q$, $\log Q$, $1/Q$, $Q^2$, $\log\sin$ of the angles, $e_3$ and the cotangents with sums over the
sides of $\log L_i$, $L_i^2$, $1/L_i$, $L_i^{1/2}$ and $L_i^3$ and with the largest side, where
$L_i=|v_i|^2/\operatorname{Im}\tau$. The best four-term form, in $Q$, $\log Q$, $\sum_i\log L_i$ and $\sum_i L_i^2$,
reaches $\chi^2$ per degree of freedom $2.58$, against $1.25$ for the fourth-degree polynomial.

\section{Discussion}

The estimates are consistent with the picture Lyons proposed: a limit law of the superbase that depends only on the
conformal class of the torus and is concentrated on the few shortest superbases. They do not prove convergence, and
the tables carry no bound on the discretization error. A natural route to a formula is through determinants of
twisted Laplacians. For cycle-rooted spanning forests the law of the homology is a ratio of such determinants, by
Forman's expansion~\cite{For93} and Kenyon's work~\cite{Ken11}, and its limit is known~\cite{KK17}. The expansion does
not apply directly here. In the line-bundle formula each cycle of a forest contributes $2-2\cos(t\theta)$, so every
Taylor coefficient in the connection counts cycle-rooted forests, and the second-order coefficient counts connected
unicyclic spanning subgraphs. The dual of a spanning tree is a connected graph with two independent cycles, which
these coefficients do not count. A formula would need a different generating object, perhaps through the period
matrices of discrete surfaces studied by Lam, Lo and Yuen~\cite{LLY26}. Whatever its form, the slope of $-\log w$ in $Q$
is $0.909$ at $\tau=i$ and $0.849$ at $\tau=2i$ on the finest meshes, so the weight is not a single exponential of the
energy.

\section*{Acknowledgements and funding}
This work was done by multiple AI agents, Claude (Anthropic) and Codex (OpenAI), connected to Math
Brain~\cite{MathBrainRepo} and working together inside the Brain framework~\cite{BrainFramework}, a shared, persistent
memory in which the agents record claims, proofs, failed attempts and methods, retrieve one another's results,
coordinate who works on what, and audit each other's work. For more on the details behind the Brain and Math Brain, see
the Math Brain paper~\cite{MathBrainPaper}, the Math Brain repository~\cite{MathBrainRepo} and the Brain framework
repository~\cite{BrainFramework}.

The research presented in this paper was supported by the European Research Council (ERC) under the European Union's
Horizon Europe research and innovation programme (grant agreement No.~101041711), by the Simons Foundation as part of
the Collaboration on the Mathematical and Scientific Foundations of Deep Learning, by Heights Labs, by Convex Nexus
Capital, by the Israel Science Foundation (grant number 2258/19), and by the Israel Science Foundation (ISF Grant
4101/25). Computational resources were provided in part by Google Cloud through its education credits program (credit
530550103).

\section*{Data availability}
The samplers, the exact enumeration, the analysis programs and the sampled class sets are available at\\
\url{https://github.com/yspennstate/ust-torus-dual-theta/tree/v1.1}.
The class sets of the $80$ tori with $10^6$ trees are attached to the v1.0 release as compressed files, and the
validation draws of Tables~\ref{tab:exact} and~\ref{tab:eight} to the v1.1 release.

The paper is part of the Math Brain project, whose public repository is given in~\cite{MathBrainRepo}.

\section*{Declaration of competing interest}
The author holds equity in Heights Labs and in Convex Nexus Capital and has a management role at Convex Nexus Capital.
These entities are included in the funding declaration above.

\section*{Declaration of generative AI and AI-assisted technologies}
Claude agents (Anthropic) wrote the programs, ran the simulations, proved Lemma~\ref{lem:struct} and prepared the
manuscript; Codex agents (OpenAI) refereed it. The question is Russell Lyons's~\cite{Ly03}, in the form given by Oded
Schramm~\cite{Sch07}. The sampler is David Wilson's algorithm~\cite{Wil96}. The determinant expansion over cycle-rooted
forests is due to Robin Forman~\cite{For93}, and the homology laws of cycle-rooted forests to Richard
Kenyon~\cite{Ken11} and to Adrien Kassel and Kenyon~\cite{KK17}. The correspondence of homological quasi-trees is due to
Wai Yeung Lam, On-Hei Solomon Lo and Chi Ho Yuen~\cite{LLY26}. The author directed the work and is responsible for the content
of the paper.

\begin{thebibliography}{99}
\bibitem{BLR25} N.~Berestycki, B.~Laslier and G.~Ray, \emph{Dimers on Riemann surfaces I: Temperleyan forests}, Ann.
Inst. Henri Poincar\'e D \textbf{12} (2025), no.~2, 363--444, \doi{10.4171/AIHPD/193}, \arxiv{1908.00832}.
\bibitem{BLR24} N.~Berestycki, B.~Laslier and G.~Ray, \emph{Dimers on Riemann surfaces II: conformal invariance and
scaling limit}, Probab. Math. Phys. \textbf{5} (2024), 961--1037, \doi{10.2140/pmp.2024.5.961}.
\bibitem{Dub15} J.~Dub\'edat, \emph{Dimers and families of Cauchy--Riemann operators I}, J. Amer. Math. Soc.
\textbf{28} (2015), 1063--1167, \doi{10.1090/jams/824}, \arxiv{1110.2808}.
\bibitem{DG15} J.~Dub\'edat and R.~Gheissari, \emph{Asymptotics of height change on toroidal Temperleyan dimer models},
J. Stat. Phys. \textbf{159} (2015), 75--100, \doi{10.1007/s10955-014-1181-x}, \arxiv{1407.6227}.
\bibitem{For93} R.~Forman, \emph{Determinants of Laplacians on graphs}, Topology \textbf{32} (1993), 35--46,
\doi{10.1016/0040-9383(93)90035-T}.
\bibitem{KK17} A.~Kassel and R.~Kenyon, \emph{Random curves on surfaces induced from the Laplacian determinant}, Ann.
Probab. \textbf{45} (2017), 932--964, \doi{10.1214/15-AOP1078}, \arxiv{1211.6974}.
\bibitem{Ken11} R.~Kenyon, \emph{Spanning forests and the vector bundle Laplacian}, Ann. Probab. \textbf{39} (2011),
1983--2017, \doi{10.1214/10-AOP596}, \arxiv{1001.4028}.
\bibitem{KSW16} R.~Kenyon, N.~Sun and D.~B. Wilson, \emph{On the asymptotics of dimers on tori}, Probab. Theory
Related Fields \textbf{166} (2016), 971--1023, \doi{10.1007/s00440-015-0687-8}, \arxiv{1310.2603}.
\bibitem{KW15} R.~Kenyon and D.~B. Wilson, \emph{Spanning trees of graphs on surfaces and the intensity of loop-erased
random walk on planar graphs}, J. Amer. Math. Soc. \textbf{28} (2015), 985--1030, \doi{10.1090/S0894-0347-2014-00819-5},
\arxiv{1107.3377}.
\bibitem{LLY26} W.~Y.~Lam, O.-H.~S.~Lo and C.~H.~Yuen, \emph{Period matrices and homological quasi-trees on discrete
Riemann surfaces}, J. Reine Angew. Math., published online 2026, \doi{10.1515/crelle-2026-0033}, \arxiv{2506.02317}.
\bibitem{Ly03} R.~Lyons, \emph{Determinantal probability measures}, Publ. Math. Inst. Hautes \'Etudes Sci.
\textbf{98} (2003), 167--212, \doi{10.1007/s10240-003-0016-0}, \arxiv{math/0204325}.
\bibitem{Sch07} O.~Schramm, \emph{Conformally invariant scaling limits: an overview and a collection of problems},
in: Proceedings of the International Congress of Mathematicians, Madrid 2006, Vol.~I, European Mathematical Society,
Z\"urich, 2007, 513--543, \doi{10.4171/022-1/20}, \arxiv{math/0602151}.
\bibitem{Wil96} D.~B. Wilson, \emph{Generating random spanning trees more quickly than the cover time}, in:
Proceedings of the 28th Annual ACM Symposium on Theory of Computing, ACM, New York, 1996, 296--303,
\doi{10.1145/237814.237880}.
\bibitem{MathBrainRepo} Y.~Shmalo, \emph{Math Brain}, public software repository, snapshot of 14 November 2026,
\url{https://github.com/yspennstate/math-brain-public/tree/8c9a6bd478c4ec4082347151b7ae16db50c917dd}.
\bibitem{MathBrainPaper} Y.~Shmalo, \emph{Math Brain: shared, auditable memory for multi-agent mathematical research},
manuscript dated 26 October 2026, Zenodo release 12 November 2026,
\url{https://github.com/yspennstate/math-brain-public/tree/main/papers/math-brain-2026-10-24},
\doi{10.5281/zenodo.24150837}.
\bibitem{BrainFramework} Y.~Shmalo, \emph{Brain Framework}, public software repository, snapshot of 13 November 2026,
\url{https://github.com/yspennstate/brain-framework-public/tree/5629e644cdd510173ae97e9aa92642dabd22dc57}.
\end{thebibliography}

\end{document}
